\documentclass[UTF8,a4paper,zihao=-4]{ctexart}

% =====================================================
% 宏包加载
% =====================================================
\usepackage{geometry}
\geometry{left=2.5cm,right=2.5cm,top=2.8cm,bottom=2.8cm}

\usepackage{fontspec}
\setmainfont{Times New Roman}

\usepackage{amsmath,amssymb,amsfonts} % 数学符号
\usepackage{mathrsfs} % 花体字母
\usepackage{graphicx} % 插入图片
\usepackage{booktabs} % 三线表
\usepackage{caption}
\usepackage{subcaption}
\usepackage{enumitem}
\usepackage{algorithm}
\usepackage{algorithmic}
\usepackage{tocloft}
\usepackage{titletoc}
\usepackage{titlesec}
\usepackage{setspace}
\usepackage{fancyhdr}
\usepackage{url}

% =====================================================
% 格式设置
% =====================================================

% --- 编号设置 ---
\numberwithin{equation}{section}
\numberwithin{figure}{section}
\numberwithin{table}{section}

% --- 目录样式 ---
\titlecontents{section}
  [0pt]
  {\addvspace{6pt}\normalfont\heiti\zihao{4}}
  {第\thecontentslabel 章\hspace{1em}}
  {}
  {\titlerule*[8pt]{.}\normalfont\zihao{-4}\fontfamily{ptm}\selectfont\contentspage}

\titlecontents{subsection}
  [2em]
  {\addvspace{2pt}\normalfont\songti\zihao{-4}}
  {\thecontentslabel\hspace{1em}}
  {}
  {\titlerule*[8pt]{.}\normalfont\zihao{-4}\fontfamily{ptm}\selectfont\contentspage}

\titlecontents{subsubsection}
  [4em]
  {\addvspace{2pt}\normalfont\songti\zihao{-4}}
  {\thecontentslabel\hspace{1em}}
  {}
  {\titlerule*[8pt]{.}\normalfont\zihao{-4}\fontfamily{ptm}\selectfont\contentspage}

% --- 章节标题样式 ---
\titleformat{\section}
  {\normalfont\heiti\zihao{-2}\centering}
  {第\thesection 章\hspace{0.5em}}
  {0pt}
  {}
  [\vspace{0.5em}]

\titleformat{\subsection}
  {\normalfont\heiti\zihao{4}}
  {\thesubsection\hspace{0.5em}}
  {0pt}
  {}
  [\vspace{0.3em}]

\titleformat{\subsubsection}
  {\normalfont\heiti\zihao{-4}}
  {\thesubsubsection\hspace{0.5em}}
  {0pt}
  {}
  [\vspace{0.2em}]

% --- 行距与页码 ---
\setstretch{1.5}
\pagestyle{plain}

% --- 重新定义引用命令 ---
% 为了配合你的手动参考文献列表
\newcommand{\mycite}[1]{\textsuperscript{\ref{#1}}}

% =====================================================
% 正文开始
% =====================================================
\begin{document}

% =====================================================
% 封面
% =====================================================
\thispagestyle{empty}

\vspace*{1cm}
\begin{flushright}
\heiti\zihao{4}
项目编号\quad\underline{20**10486***}
\end{flushright}

\vspace{2cm}

\begin{center}
    \songti\zihao{1}
    \textbf{武汉大学大学生创新创业训练}
    
    \textbf{计划项目中期报告}

\vspace{2cm}

\heiti\zihao{2}
\textbf{项目名称}

\vspace{2cm}

\songti\zihao{-3}
\begin{tabular}{rl}
院（系）名称： &  \\
专业名称： &  \\
学生姓名： &  \\
指导教师： &  \quad 教授 \\
\end{tabular}

\vfill

\songti\zihao{-2}
二〇二六年二月

\end{center}
\newpage

% =====================================================
% 英文扉页
% =====================================================
\thispagestyle{empty}
\begin{center}

\fontfamily{ptm}\selectfont\zihao{2}INTERIM REPORT OF PLANNING \\PROJECT OF INNOVATION AND\\
ENTREPRENEURSHIP TRAINING OF \\UNDERGRADUATES OF WUHAN \\UNIVERSITY
\end{center}

\vspace{1.5cm}
\begin{center}
    \usefont{T1}{ptm}{b}{n}\zihao{2}
    
    \textbf{Project Title}
\end{center}

\vspace{2cm}
\begin{center}
    \bfseries\fontfamily{ptm}\selectfont\zihao{4}
\begin{tabular}{rl}
College: &  \\
Subject: &  \\
Name: &  \\
Director: & \\
\end{tabular}
\end{center}

\vfill

\begin{center}
    \bfseries\fontfamily{ptm}\selectfont\zihao{-2}
    February 2026
\end{center}

\newpage

% =====================================================
% 声明
% =====================================================
\thispagestyle{empty}

\section*{\songti\zihao{2} \textbf{郑\ 重\ 声\ 明}}
\vspace{2cm}

\songti\zihao{4}
\linespread{1.5}\selectfont
本项目组呈交的中期报告，是在导师的指导下，独立进行研究工作所取得的成果，所有数据、图片资料真实可靠。尽我们所知，除文中已经注明引用的内容外，本报告的研究成果不包含他人享有著作权的内容。对本报告所涉及的研究工作做出贡献的其他个人和集体，均已在文中以明确的方式标明。本报告的知识产权归属于培养单位。

\vspace{3cm}

\begin{center}
项目组签名：\underline{\hspace{4cm}} \qquad 日期：\underline{\hspace{4cm}}
\end{center}

\newpage

% =====================================================
% 中文摘要
% =====================================================
\cleardoublepage
\pagenumbering{Roman}
\setcounter{page}{1}


\begin{center}
\heiti\zihao{-2} 摘\qquad 要
\end{center}
\addcontentsline{toc}{section}{摘要}

\songti\zihao{-4}
在此处填写摘要，建议用 300--500 字概括：项目背景、研究目标、目前已经完成的工作、阶段性结果以及下一步计划

\vspace{0.5cm}
\noindent\heiti\zihao{-4}关键词：\songti\zihao{-4}
关键词1；关键词2；关键词3；关键词4；关键词5

\newpage

% =====================================================
% 英文摘要
% =====================================================

\begin{center}
{\usefont{T1}{ptm}{b}{n}\zihao{-2} ABSTRACT}
\end{center}
\addcontentsline{toc}{section}{ABSTRACT}

\fontfamily{ptm}\selectfont\zihao{-4}
Write an English abstract corresponding to the Chinese abstract

\vspace{0.5cm}
\noindent{\usefont{T1}{ptm}{b}{n}\zihao{-4}Key words:}
\zihao{-4}Keyword 1; Keyword 2; Keyword 3; Keyword 4; Keyword 5

\newpage

% =====================================================
% 目录
% =====================================================
\clearpage
\begin{center}
    \heiti\zihao{-2}\textbf{目\hspace{1em}录}
\end{center}
\vspace{2em}
\makeatletter
\@starttoc{toc}
\makeatother
\newpage

% =====================================================
% 正文
% =====================================================
\cleardoublepage
\pagenumbering{arabic}
\setcounter{page}{1}

% =====================================================
% 第1章 绪论
% =====================================================
\section{绪论}

\subsection{研究背景}
\songti\zihao{-4}
介绍项目所属领域、现实或学术背景、当前研究现状，以及为什么需要开展本项目

\subsection{研究意义}
\songti\zihao{-4}
说明项目的理论意义、应用价值或工程价值

\subsection{研究内容}
\songti\zihao{-4}
目前各项任务推进到了什么程度

\subsection{研究方法}
\songti\zihao{-4}
简要给出整体研究方法和技术路线



% =====================================================
% 第2章 模型与算法
% =====================================================
\newpage
\section{模型与算法}

\subsection{}
\songti\zihao{-4}


如果需要给出数学模型，可以使用如下格式：
\begin{equation}
    \min_{\theta} \; \mathcal{L}(\theta)
    = \frac{1}{N}\sum_{i=1}^{N}\ell\bigl(f_{\theta}(x_i),y_i\bigr)
    + \lambda R(\theta),
    \label{eq:example_objective}
\end{equation}
其中，$\theta$ 表示模型参数，$\mathcal{L}$ 表示目标函数，$R(\theta)$ 表示正则项。正文中可通过式~\eqref{eq:example_objective} 引用该公式。



\subsection{}
\songti\zihao{-4}
算法伪代码示例如算法~\ref{alg:example} 所示。

\begin{algorithm}[htbp]
\renewcommand{\algorithmicrequire}{\textbf{输入：}}
\renewcommand{\algorithmicensure}{\textbf{输出：}}
\caption{示例算法}
\label{alg:example}
\begin{algorithmic}[1]
    \REQUIRE 输入数据 $X$，最大迭代次数 $T$
    \ENSURE 输出结果 $Y$
    \STATE 初始化模型参数
    \FOR{$t=1,2,\dots,T$}
        \STATE 根据当前参数计算目标函数
        \STATE 更新模型参数
    \ENDFOR
    \RETURN $Y$
\end{algorithmic}
\end{algorithm}


% =====================================================
% 第3章 数值实验与结果分析
% =====================================================
\newpage
\section{数值实验与结果分析}
要有实验设计的细节和实验结果的分析


\subsection{数值实验1}
\songti\zihao{-4}



\newpage

\subsection{数值实验2}
\songti\zihao{-4}


% =====================================================
% 第4章 项目进度与后续计划
% =====================================================
\newpage
\section{项目进度与后续计划}

\subsection{项目完成情况}
\songti\zihao{-4}



\subsection{目前存在的问题}
\songti\zihao{-4}



\subsection{下一阶段计划}
\songti\zihao{-4}

% =====================================================
% 参考文献
% =====================================================
\newpage
\section*{\heiti\zihao{-2} 参考文献}
\addcontentsline{toc}{section}{参考文献}

\songti\zihao{-4}
\begin{enumerate}[
    label={[\arabic*]}, 
    leftmargin=3.5em, 
    itemsep=4pt, 
    align=left
]
    \item \label{cite:DeepRitz} E Weinan, Yu Bing. The Deep Ritz Method: A Deep Learning-Based Numerical Algorithm for Solving Variational Problems [J]. Communications in Mathematics and Statistics, 2018, 6(1): 1-12.
\end{enumerate}

\end{document}